\section{Anexo de matriz de riesgos integrada}

La matriz consolidada identifica riesgos técnicos, legales, financieros, de proyecto y éticos asociados a los entregables D7--D14. La escala empleada es de 1 a 3 para probabilidad e impacto; la prioridad corresponde al producto de ambas dimensiones \parencite{vicente-riesgos}.

\begin{table}[H]
  \centering
  \caption{Matriz de riesgos priorizada: R01--R05}\label{tab:anexo-riesgos-integrados-a}
  \scriptsize
  \begin{tabular}{@{}l l c l l@{}}
    \toprule
    ID & Categoría & P$\times$I & Causa y consecuencia & Mitigación y responsable \\
    \midrule
    R01 & Proyecto / financiero & 2$\times$3 = 6 (Crítico) & \parbox[t]{0.31\textwidth}{Mayor complejidad técnica incrementa horas de RRHH e impacta el 69,7\% del TCO base.} & \parbox[t]{0.29\textwidth}{Validar horas en EDT/Gantt, registrar cambios de alcance y monitorear semanalmente. \textit{Fabián / Ignacio.}} \\
    R02 & Legal & 2$\times$3 = 6 (Crítico) & \parbox[t]{0.31\textwidth}{La aprobación para integrar datos de salud en GCP puede retrasar el cronograma de diez semanas.} & \parbox[t]{0.29\textwidth}{Aprobación formal temprana y pseudonimización estricta antes de cargar datos. \textit{Vicente.}} \\
    R03 & Financiero & 3$\times$2 = 6 (Crítico) & \parbox[t]{0.31\textwidth}{La variación USD/CLP puede desviar el OPEX planificado.} & \parbox[t]{0.29\textwidth}{Contingencia de 10\%, alertas de gasto y revisión mensual. \textit{Ignacio.}} \\
    R04 & Financiero & 2$\times$2 = 4 (Medio) & \parbox[t]{0.31\textwidth}{Egress, logs o servicios administrados no previstos pueden producir cobros adicionales.} & \parbox[t]{0.29\textwidth}{Etiquetar costos, asignar presupuesto por servicio y aprobar activaciones. \textit{Ignacio.}} \\
    R05 & Técnico / proyecto & 1$\times$2 = 2 (Bajo) & \parbox[t]{0.31\textwidth}{Descalces de la ingesta batch pueden generar inconsistencia de datos al cierre de jornada.} & \parbox[t]{0.29\textwidth}{Monitorear pipeline y cumplir RTO menor a 8 horas. \textit{Lucas / Fabián.}} \\
    \bottomrule
  \end{tabular}
\end{table}

\begin{table}[H]
  \centering
  \caption{Matriz de riesgos priorizada: R06--R10}\label{tab:anexo-riesgos-integrados-b}
  \scriptsize
  \begin{tabular}{@{}l l c l l@{}}
    \toprule
    ID & Categoría & P$\times$I & Causa y consecuencia & Mitigación y responsable \\
    \midrule
    R06 & Financiero / legal & 2$\times$1 = 2 (Bajo) & \parbox[t]{0.31\textwidth}{El tratamiento tributario de servicios Cloud podría diferir del 19\% estimado.} & \parbox[t]{0.29\textwidth}{Validar el tratamiento con administración institucional antes de contratar. \textit{Ignacio / Vicente.}} \\
    R07 & Técnico & 2$\times$3 = 6 (Crítico) & \parbox[t]{0.31\textwidth}{Si no se aprueba nube pura, la sincronización entre entorno local y GCP puede fallar.} & \parbox[t]{0.29\textwidth}{Probar sincronización tempranamente y asignar un responsable técnico único. \textit{Lucas.}} \\
    R08 & Legal & 1$\times$3 = 3 (Medio) & \parbox[t]{0.31\textwidth}{Una brecha de seguridad podría exponer registros clínicos identificables.} & \parbox[t]{0.29\textwidth}{Aplicar IAM de mínimo privilegio, Secret Manager y pseudonimización previa. \textit{Vicente.}} \\
    R09 & Técnico / ético & 2$\times$2 = 4 (Medio) & \parbox[t]{0.31\textwidth}{La deriva de datos puede degradar la precisión durante los 36 meses de operación.} & \parbox[t]{0.29\textwidth}{Realizar seis ventanas semestrales de reentrenamiento y aprobación médica obligatoria. \textit{César / Lucas.}} \\
    R10 & Ético & 2$\times$3 = 6 (Crítico) & \parbox[t]{0.31\textwidth}{La cobertura del dataset académico deduplicado es desigual por sexo y limitada en edades extremas e intersecciones; además, su generalización a un hospital chileno es incierta.} & \parbox[t]{0.29\textwidth}{Evaluar desempeño por subgrupos con N suficiente, validar localmente, monitorear deriva y mantener supervisión humana. SHAP apoya explicabilidad, no prueba equidad. \textit{César; riesgo residual por definir.}} \\
    \bottomrule
  \end{tabular}
\end{table}

\noindent\textit{Nota:} el informe adopta P$\times$I = 2$\times$3 = 6 (Crítico) como nivel propuesto/final de R10. El riesgo residual no está definido y no se infiere en esta integración.

\subsection{Criterios de aceptación de riesgos}

La matriz satisface los criterios definidos por su fuente: diferencia riesgos técnicos, legales, financieros y éticos; mantiene una escala consistente de probabilidad e impacto; y vincula cada riesgo con una mitigación y un responsable. Cloud/GCP es la arquitectura seleccionada para el piloto, condicionada a los gates éticos, legales, técnicos, de seguridad, de riesgos y de validación local. Los riesgos críticos R01, R02, R03, R07 y R10 deben revisarse en los gates correspondientes antes de avanzar al piloto o a una nueva versión del modelo. El riesgo residual de R10 permanece sin definir.

\subsection{Trazabilidad financiera histórica}

El TCO oficial Cloud para este informe es \$10.815.487 CLP a 36 meses, según la evaluación financiera D7--D11. El valor de \$11.866.338 presente en una versión previa de la matriz de riesgos se conserva únicamente como antecedente histórico de alineación; no es un valor vigente ni condiciona el presupuesto del piloto.
